% Generated from common/cover-notes.json by tools/series_cover.py.
\backmatter
\cleardoublepage
\chapter*{About the Cover}
\addcontentsline{toc}{chapter}{About the Cover}
\noindent The cover draws regular heptagons in the Poincaré disc, with three meeting at each vertex. They have the same hyperbolic size, although the picture makes them look smaller near the boundary. The boundary circle itself is infinitely far away. Poincaré used this disc model in his 1882 work on Fuchsian groups. Identifying tiles under an infinite discrete group of hyperbolic isometries gives Klein’s quartic, a surface with three handles and a finite heptagon tiling. Chapter 9 connects such identifications with matrices and geodesic motion. The book's own TikZ cover generator computed the drawing.
